\chapter{هنرِ حلِ مسئله و گام‌های بعدی}%
\label{ch:problem-solving}

به فصلِ پایانی رسیدیم. دیگر همهٔ ابزارهای پایه را دارید؛ آنچه می‌ماند
مهارت است: اینکه وقتی با مسئله‌ای تازه و ناآشنا روبه‌رو شدید، از کجا شروع
کنید و وقتی گیر کردید چه کنید. این فصل چکیدهٔ تجربهٔ نسل‌ها برنامه‌نویس
است، به‌اضافهٔ یک پروژهٔ کاملِ پایانی که همه‌چیز را کنارِ هم می‌نشاند.

\section{دستورالعملِ رویارویی با مسئلهٔ تازه}

\begin{enumerate}
  \item \textbf{مسئله را بفهمید.} پیش از هر چیز، با کلماتِ خودتان
        بازگویش کنید: ورودی دقیقاً چیست؟ خروجی دقیقاً چیست؟ برای دو سه
        ورودیِ نمونه، خروجی را \emph{با دست} حساب کنید. اگر نتوانید با
        دست حساب کنید، هرگز نمی‌توانید به رایانه یادش بدهید.
  \item \textbf{مسئله را کوچک کنید.} مسئلهٔ بزرگ را به زیرمسئله‌های
        مستقل بشکنید و به هر کدام نامی بدهید؛ این نام‌ها معمولاً همان
        کارکردهای برنامهٔ شما می‌شوند. اگر زیرمسئله‌ای هنوز بزرگ بود، باز
        هم بشکنیدش.
  \item \textbf{از حالتِ ساده شروع کنید.} می‌خواهید برنامه با هر تعداد
        عدد کار کند؟ اول با سه عددِ ثابت راهش بیندازید. نسخهٔ ساده که
        کار کرد، عمومی‌اش کنید.
  \item \textbf{شبه‌کد بنویسید.} پیش از تایپ، راهِ‌حل را در چند خط
        شبه‌کد بیاورید و یک اجرای دستی کنید؛ همان عادتِ فصلِ
        \ref{ch:algorithms}.
  \item \textbf{بنویسید، اجرا کنید، تکرار کنید.} همه‌چیز را یک‌جا ننویسید.
        چند خط بنویسید، اجرا کنید، مطمئن شوید، بعد ادامه دهید. هر چه
        فاصلهٔ دو اجرا کوتاه‌تر باشد، پیدا کردنِ اشکال آسان‌تر است، چون
        مقصر میانِ همان چند خطِ آخر است.
\end{enumerate}

\section{وقتی برنامه درست کار نمی‌کند}

خواهد شد؛ به همهٔ ما می‌شود، هر روز. به اشکالِ برنامه \textbf{باگ}
می‌گویند و به یافتن و رفعش \textbf{اشکال‌زدایی}. چند اصل:

\begin{itemize}
  \item \textbf{پیامِ خطا را بخوانید.} جدی و کامل، نه فقط نگاهی از سرِ
        ترس. شمارهٔ خط و توضیحِ خطا نیمی از راه است.
  \item \textbf{چاپِ ردیابی بگذارید.} ساده‌ترین و کهن‌ترین ابزارِ
        اشکال‌زدایی همان \code{چاپ} است: وسطِ برنامه مقدارِ متغیرهای
        مشکوک را چاپ کنید و ببینید از کجا به بعد، واقعیت از انتظارِ شما
        جدا می‌شود. همان‌جا لانهٔ باگ است.
  \item \textbf{به حالت‌های مرزی بدبین باشید.} صفر، منفی، آرایهٔ خالی،
        دو مقدارِ برابر، اولین و آخرین خانه. در تمرین‌های این کتاب بارها
        دیدید که باگ‌ها همین گوشه‌ها قایم می‌شوند.
  \item \textbf{برای یک نفر توضیحش بدهید.} خط‌به‌خط، بلندبلند؛ حتی برای
        یک عروسک! جایی که وسطِ توضیح مکث می‌کنید و می‌گویید «اِ، صبر
        کن\ldots»، همان‌جاست. این ترفند آن‌قدر کار می‌کند که نامِ جهانی
        دارد: اشکال‌زداییِ اردک‌پلاستیکی.
\end{itemize}

\section{پروژهٔ پایانی: کارنامه‌ساز}

مسئله: نمره‌های یک کلاس را داریم. می‌خواهیم گزارشی بسازیم شاملِ میانگین،
بالاترین و پایین‌ترین نمره، تعدادِ قبولی‌ها، و فهرستِ مرتب‌شدهٔ نمره‌ها.

طبقِ دستورالعمل، مسئله را به زیرمسئله‌ها می‌شکنیم و خوشبختانه همه را در
این کتاب ساخته‌ایم! فقط باید کنارِ هم بنشینند:

\begin{salam}
کارکرد میانگین(نمرات : صحیح[:]) : اعشار:
    گذرا مجموع : صحیح = 0
    هر نمره در نمرات:
        مجموع = مجموع + نمره
    پایان
    بازگشت مجموع * 1.0 / len(نمرات)
پایان

کارکرد بیشینه آرایه(نمرات : صحیح[:]) : صحیح:
    گذرا بزرگترین : صحیح = نمرات[0]
    هر نمره در نمرات:
        اگر نمره > بزرگترین:
            بزرگترین = نمره
        پایان
    پایان
    بازگشت بزرگترین
پایان

کارکرد کمینه آرایه(نمرات : صحیح[:]) : صحیح:
    گذرا کوچکترین : صحیح = نمرات[0]
    هر نمره در نمرات:
        اگر نمره < کوچکترین:
            کوچکترین = نمره
        پایان
    پایان
    بازگشت کوچکترین
پایان

کارکرد شمارش قبولی(نمرات : صحیح[:]) : صحیح:
    گذرا تعداد : صحیح = 0
    هر نمره در نمرات:
        اگر نمره >= 10:
            تعداد = تعداد + 1
        پایان
    پایان
    بازگشت تعداد
پایان

کارکرد مرتب سازی حبابی(نمرات : صحیح[:]):
    ن : صحیح = len(نمرات)
    چرخه ن - 1:
        گذرا i : صحیح = 0
        تاهنگام i < ن - 1:
            اگر نمرات[i] > نمرات[i + 1]:
                موقت : صحیح = نمرات[i]
                نمرات[i] = نمرات[i + 1]
                نمرات[i + 1] = موقت
            پایان
            i = i + 1
        پایان
    پایان
پایان

کارکرد آغازین:
    گذرا نمرات : صحیح[8] = [14, 9, 18, 11, 7, 20, 15, 12]
    چاپ "===== کارنامه کلاس ====="
    چاپ "تعداد دانش‌آموزان:", len(نمرات)
    چاپ "میانگین:", میانگین(نمرات)
    چاپ "بالاترین نمره:", بیشینه آرایه(نمرات)
    چاپ "پایین‌ترین نمره:", کمینه آرایه(نمرات)
    چاپ "تعداد قبولی:", شمارش قبولی(نمرات)
    مرتب سازی حبابی(نمرات)
    بنویس "نمره‌ها از کم به زیاد: "
    هر نمره در نمرات:
        بنویس نمره
        بنویس " "
    پایان
    چاپ ""
پایان
\end{salam}

\begin{progoutfa}
===== کارنامه کلاس =====
تعداد دانش‌آموزان: 8
میانگین: 13.25
بالاترین نمره: 20
پایین‌ترین نمره: 7
تعداد قبولی: 6
نمره‌ها از کم به زیاد: 7 9 11 12 14 15 18 20
\end{progoutfa}

به دو ریزه‌کاری دقت کنید. یکی در کارکردِ \code{میانگین}: ضرب در \code{1.0}
مجموعِ صحیح را اعشاری می‌کند تا گرفتارِ دامِ تقسیمِ صحیحِ فصلِ
\ref{ch:operators} نشویم و ۱۳٫۲۵ بگیریم، نه ۱۳. دیگری در کارکردِ آغازین:
کارکردِ آغازین خودش تقریباً هیچ کاری نمی‌کند؛ فقط رهبرِ ارکستر است و هر قطعه
را کارکردی می‌نوازد که جداگانه نوشته و آزموده شده است. برنامهٔ خوب همین
شکلی است.

\section{از اینجا به کجا؟}

شما اکنون چیزی دارید که هیچ‌کس نمی‌تواند از شما بگیرد: \emph{سوادِ
الگوریتمی}. متغیر و شرط و حلقه و کارکرد و آرایه در هر زبانِ برنامه‌نویسیِ
دیگری هم (با لباسی دیگر) همین است که آموختید. برای ادامهٔ راه:

\begin{itemize}
  \item \textbf{بنویسید، زیاد بنویسید.} پروژه‌های کوچکِ خودتان را
        بسازید: بازیِ حدسِ عدد، آزمونِ جدول‌ضرب، دفترِ نمره‌ها. مسئله
        از زندگیِ خودتان بهتر از هر تمرینِ کتابی است.
  \item \textbf{تمرین‌های حل‌نشده را حل کنید.} اگر تمرینی را رد
        کرده‌اید، برگردید؛ به‌خصوص چالشی‌ها را.
  \item \textbf{بازنویسی کنید.} برنامه‌های همین کتاب را از حفظ و به
        سلیقهٔ خودتان دوباره بنویسید؛ چیزهایی یاد می‌گیرید که بارِ اول
        دیده نمی‌شوند.
  \item \textbf{سراغِ جلدِ بعد بروید.} زبانِ سلام بسیار بیش از این‌هاست:
        رشته‌ها و پردازشِ متن، ساختارها، کار با فایل‌ها و کتابخانهٔ
        استاندارد. کتابِ «راهنمای کاملِ آموزشیِ زبانِ سلام» ادامهٔ
        طبیعیِ همین کتاب است.
\end{itemize}

\begin{center}
\vspace{0.5cm}
\emph{روزی که این کتاب را باز کردید، رایانه برایتان جعبه‌ای جادویی بود.\\
امروز می‌دانید جادویی در کار نیست: فکرِ روشن است و گام‌های کوچک.\\
این یعنی شما دیگر برنامه‌نویسید. به دنیای ما خوش آمدید!}
\vspace{0.5cm}
\end{center}

\section{تمرین‌های پایانی}

\exercise به کارنامه‌ساز قابلیتِ تازه بیفزایید: چاپِ «فاصلهٔ نمره‌ها»،
یعنی تفاضلِ بالاترین و پایین‌ترین نمره (بدونِ نوشتنِ حلقهٔ تازه؛ از
کارکردهای موجود استفاده کنید).

\exercise کارکردی بنویسید که تعدادِ نمره‌های \emph{تک‌رقمی} را برگرداند و
به گزارش اضافه کنید.

\exercise بازیِ ذهنی: برنامه‌ای بنویسید که همهٔ اعدادِ سه‌رقمی‌ای را
پیدا کند که با وارونهٔ خود برابرند و اول هم هستند (کارکردهای \code{وارونه}
و \code{اول است} را از فصل‌های پیش دارید).

\exercise (چالشیِ بزرگ) مرتب‌سازیِ انتخابی: روشِ دیگری برای مرتب‌سازی
چنین است: کوچک‌ترین عضوِ آرایه را پیدا کن و با خانهٔ اول جابه‌جا کن؛
سپس کوچک‌ترینِ باقی‌مانده را با خانهٔ دوم؛ و همین‌طور تا آخر. آن را
پیاده کنید و خروجی‌اش را با مرتب‌سازیِ حبابی مقایسه کنید.
